\documentclass[11pt,oneside]{amsbook}

\usepackage[T1]{fontenc}
\usepackage[utf8]{inputenc}
\usepackage{lmodern}
\usepackage{microtype}
\usepackage{amsmath,amssymb,mathtools}
\usepackage[
  a4paper,
  inner=31mm,
  outer=31mm,
  top=28mm,
  bottom=30mm,
  marginparwidth=22mm,
  marginparsep=5mm
]{geometry}
\usepackage[hidelinks,pdfencoding=auto]{hyperref}
\usepackage{bookmark}

\newtheorem{proposition}{Proposition}[section]
\newtheorem{theorem}[proposition]{Theorem}
\newtheorem{corollary}[proposition]{Corollary}
\newtheorem{definition}[proposition]{Definition}
\newtheorem{corollaries}[proposition]{Corollaries}

\hypersetup{
  pdftitle={SGA 1 English source-synchronization workpass: Exposé I opening through I.6},
  pdfauthor={Working translation from the French arXiv authority},
  pdfsubject={Bounded source-synchronization English working translation},
  pdfkeywords={SGA 1, étale morphisms, source synchronization, working translation}
}

\renewcommand{\thechapter}{\Roman{chapter}}
\renewcommand{\thesection}{\thechapter.\arabic{section}}
\setcounter{secnumdepth}{1}
\setcounter{tocdepth}{1}

\begin{document}

\frontmatter
\begin{titlepage}
  \centering
  \vspace*{35mm}
  {\Huge\bfseries \textit{Étale Coverings and the Fundamental Group}\par}
  \vspace{3mm}
  {\LARGE (SGA~1)\par}
  \vspace{20mm}
  {\Large English source-synchronization workpass\par}
  \vspace{3mm}
  {\Large Exposé~I opening through \S\,I.6\par}
  \vfill
  {\large Machine-assisted bounded working translation\par}
  \vspace{5mm}
  {\large 19 July 2026\par}
\end{titlepage}

\chapter*{Source and status note}

This bounded workpass translates the opening of Exposé~I through \S\,I.6 directly
from the French \TeX{} source in arXiv \texttt{math/0206203v2}, with archive
SHA-256:
\begin{center}
  \small\ttfamily
  0A33C9C06908705A2525690FAAE02F6F\\
  07980A4AA069A8B6CFF9B1D9BC39ACD3
\end{center}
The current translated source boundary is
\texttt{smf\_doc-math\_3\_01.tex}, lines 556--1216; line 1217 begins \S\,I.7
and is outside the current workpass.  The opening through \S\,I.5 is preserved as prior source-reviewed history.
The \S\,I.6 addition extends the current public bounded checkpoint and is the
current source-, formula-, structure-, terminology-, build-, and render-review
unit.

Recovered English materials associated with jcreinhold, jmoellermath, and
Tim Hosgood were consulted only as target-language comparison controls.
They are not source authority, and no agreement among them was treated as
proof of source alignment.  The two Hosgood-related repositories belong to
one derivative lineage, not two independent witnesses.

This is a source-audited working translation in bounded units.  It is not
a complete SGA~1 translation, a critical edition, or a certification of the
mathematics.  This checkpoint does not determine redistribution rights or grant a blanket
license over the French source or new English translation.  The French source
header asks modified versions to identify themselves, which this workpass does
explicitly.

\tableofcontents

\mainmatter
\setcounter{chapter}{0}
\chapter{Étale morphisms}
\label{I}
\marginpar{\raggedright\footnotesize French source p.~1}

To simplify the exposition, we assume that all preschemes under
consideration are locally noetherian, at least after no.~I.2.

\section{Notions of differential calculus}
\label{I.1}

Let \(X\) be a prescheme over \(Y\), and let \(\Delta_{X/Y}\), or simply
\(\Delta\),\label{indnot:ab} be the diagonal morphism
\(X\to X\times_Y X\).  It is an
immersion, hence a closed immersion of \(X\) into an open subset \(V\) of
\(X\times_Y X\).  Let \(\mathcal I_X\) be the ideal of the closed
subprescheme corresponding to the diagonal in \(V\) (N.B. if one wishes to
proceed intrinsically, without assuming that \(X\) is separated over
\(Y\)---a hypothesis that would be farcical---one should consider the
set-theoretic inverse image of \(\mathcal O_{X\times X}\) in \(X\), and
denote by \(\mathcal I_X\) the augmentation ideal in the latter\ldots).
The sheaf \(\mathcal I_X/\mathcal I_X^2\) may be regarded as a
quasi-coherent sheaf on \(X\); we denote it by \(\Omega^1_{X/Y}\).  It is
\label{indnot:ac}
of finite type if \(X\to Y\) is of finite type.  It behaves well under base
change \(Y'\to Y\).  We also introduce the sheaves
\(\mathcal O_{X\times_Y X}/\mathcal I_X^{n+1}=\mathcal P^n_{X/Y}\);\label{indnot:ad} these
are sheaves of \emph{rings} on \(X\), making \(X\) into a prescheme that
may be denoted by \(\Delta^n_{X/Y}\)\label{indnot:ae} and called the \emph{\(n\)-th
infinitesimal neighborhood of \(X/Y\)}.  The resulting sorites is entirely
trivial, although rather long\footnote{Cf.\ EGA IV~16.3.}; it would be
prudent to speak of it only when there is something useful to say about it,
in connection with smooth morphisms.

\section{Quasi-finite morphisms}
\label{I.2}

\begin{proposition}
\label{I.2.1}
Let \(A\to B\) be a local homomorphism (N.B. the rings are now
noetherian), and let \(\mathfrak m\) be the maximal ideal of~\(A\).  The
following conditions are equivalent:
\begin{enumerate}
\item[(i)] \(B/\mathfrak mB\) is finite-dimensional over
  \(k=A/\mathfrak m\).
\item[(ii)] \(\mathfrak mB\) is an ideal of definition, and
  \(B/\mathfrak r(B)=\kappa(B)\) is an extension of
  \(k=\kappa(A)\).
\item[(iii)] The completion \(\widehat B\) is finite over the completion
  \(\widehat A\) of~\(A\).
\end{enumerate}
\end{proposition}

One then says\marginpar{\raggedright\footnotesize French source p.~2} that
\(B\) \emph{is quasi-finite} over~\(A\).  A morphism
\(f\colon X\to Y\) is said to be quasi-finite at~\(x\) (or the
\(Y\)-prescheme \(f\) is said to be quasi-finite at~\(x\)) if
\(\mathcal O_x\) is quasi-finite over \(\mathcal O_{f(x)}\).  This is also
equivalent to saying that \(x\) is \emph{isolated in its fiber}
\(f^{-1}(f(x))\).\footnote{The original printing (Exposé~I, p.~2) and both
French arXiv renderings print \(f^{-1}(x)\).  Since \(x\in X\) for
\(f\colon X\to Y\), the type-correct fiber through \(x\) is
\(f^{-1}(f(x))\), which is used here with the source defect explicitly
disclosed.} A morphism is said to be quasi-finite if it is so at
every point.\footnote{In EGA~II~6.2.3 one assumes in addition that \(f\) is
of finite type.}

\begin{corollary}
\label{I.2.2}
If \(A\) is complete, quasi-finite is equivalent to finite.
\end{corollary}

One could give the usual chain of implications \textup{(i), (ii), (iii),
(iv), (v)} for quasi-finite morphisms, but this does not seem indispensable
here.

% SGA 1, Expose I, section I.3.
% French authority lines 653--799; source-audited bounded unit.

\section{Unramified or net morphisms}
\label{I.3}

\begin{proposition}
\label{I.3.1}
Let \(f\colon X\to Y\) be a morphism of finite type, let \(x\in X\), and
put \(y=f(x)\).  The following conditions are equivalent:
\begin{enumerate}
\item[(i)] 
  \(\mathcal O_x/\mathfrak m_y\mathcal O_x\) is a finite separable
  extension of~\(\kappa(y)\).
\item[(ii)] \(\Omega^1_{X/Y}\) vanishes at~\(x\).
\item[(iii)] The diagonal morphism \(\Delta_{X/Y}\) is an open immersion
  in a neighborhood of~\(x\).
\end{enumerate}
\end{proposition}

For the implication \textup{(i)}\(\Rightarrow\)\textup{(ii)}, Nakayama's
lemma immediately reduces us to the case
\(Y=\operatorname{Spec}(k)\), \(X=\operatorname{Spec}(k')\), where the
assertion is well known and, moreover, immediate from the definition of
separability.  The implication \textup{(ii)}\(\Rightarrow\)\textup{(iii)}
follows from a pleasant and easy characterization of open immersions using
Krull.  For \textup{(iii)}\(\Rightarrow\)\textup{(i)}, one is
again reduced to the case \(Y=\operatorname{Spec}(k)\), with the diagonal
morphism an open immersion everywhere.  One must then prove that \(X\) is
finite, with a separable coordinate ring over~\(k\); for this one reduces to
the case where \(k\) is algebraically closed.  But then every closed point of
\(X\) is isolated (it is precisely the inverse image of the diagonal under
the morphism \(X\to X\times_k X\) defined by~\(x\)), and hence \(X\) is
finite.  One may therefore suppose that \(X\) consists of a single point,
with ring~\(A\).  Thus \(A\otimes_kA\to A\) is an isomorphism, whence
\(A=k\), as required.

\begin{definition}
\label{I.3.2}
\begin{enumerate}
\item[(a)] One then says that \(f\) is \emph{net}, or equivalently
  \emph{unramified}, at~\(x\), or that \(X\) is net, or unramified, at
  \(x\) over~\(Y\).
\item[(b)] Let \(A\to B\) be a local homomorphism.  It is said to be
  \emph{net}, or \emph{unramified}, and \(B\) is said to be a net, or
  unramified, local algebra over~\(A\), if
  \(B/\mathfrak r(A)B\) is a finite separable extension of
  \(A/\mathfrak r(A)\); that is, if
  \(\mathfrak r(A)B=\mathfrak r(B)\) and \(\kappa(B)\) is a separable
  extension of~\(\kappa(A)\).\footnote{Cf.\ the regrets in
  Expos\'e~III, no.~1.2.}
\end{enumerate}
\end{definition}

\paragraph{Remarks.}
The fact\marginpar{\raggedright\footnotesize French source p.~3} that
\(B\) is unramified over~\(A\) can already be checked on the completions of
\(A\) and~\(B\).  Unramified implies quasi-finite.

\begin{corollary}
\label{I.3.3}
The set of points at which \(f\) is unramified is open.
\end{corollary}

\begin{corollary}
\label{I.3.4}
Let \(X'\) and \(X\) be two preschemes of finite type over~\(Y\), and let
\(g\colon X'\to X\) be a \(Y\)-morphism.  If \(X\) is unramified over
\(Y\), the graph morphism

\[
  \Gamma_g\colon X'\longrightarrow X'\times_Y X
\]

is an open immersion.\footnote{The original printing (Expos\'e~I, p.~3) and
the French TeX give \(X\times_Y X\) as the graph target and later write
\(\operatorname{id}_{X'}\) in the pullback map.  The typed fiber products
force \(X'\times_Y X\) and \(\operatorname{id}_X\), respectively; those
type-correct readings are used here with both source defects disclosed.}
\end{corollary}

Indeed, it is the inverse image of the diagonal morphism
\(X\to X\times_Y X\) under
\[
g\times_Y\operatorname{id}_{X}\colon
X'\times_Y X\longrightarrow X\times_Y X.
\]

One may also introduce the annihilator ideal
\(\mathfrak d_{X/Y}\)\label{indnot:af} of \(\Omega^1_{X/Y}\), called the
\emph{different ideal} of~\(X/Y\).  It defines a closed subprescheme of
\(X\) whose underlying set is the set of points at which \(X/Y\) is
ramified, that is, not unramified.

\begin{proposition}
\label{I.3.5}
\begin{enumerate}
\item[(i)] An immersion is unramified.
\item[(ii)] The composite of two unramified morphisms is unramified.
\item[(iii)] A base change of an unramified morphism is unramified.
\end{enumerate}
\end{proposition}

This may be seen equally well from criterion \textup{(ii)} or criterion
\textup{(iii)} (the latter seems more amusing to me).  One can of course
also be more precise by giving pointwise statements; this is more general
only in appearance (except in the setting of Definition~\ref{I.3.2}(b)), and
tedious.  As usual, one obtains the following corollaries.

\begin{corollaries}
\label{I.3.6}
\begin{enumerate}
\item[(iv)] The Cartesian product of two unramified morphisms is unramified.
\item[(v)] If \(gf\) is unramified, then \(f\) is unramified.
\item[(vi)] If \(f\) is unramified, then \(f_{\mathrm{red}}\) is
  unramified.
\end{enumerate}
\end{corollaries}

\begin{proposition}
\label{I.3.7}
Let \(A\to B\) be a local homomorphism.  Suppose that the residue-field
extension \(\kappa(B)/\kappa(A)\) is trivial, or that \(\kappa(A)\) is
algebraically closed.  Then \(B/A\) is unramified if and only if
\(\widehat B\), as an \(\widehat A\)-algebra, is a quotient of
\(\widehat A\).
\end{proposition}

\paragraph{Remarks.}
\begin{itemize}
\item When the residue-field extension is not assumed to be trivial, one may
  reduce to that case by making a suitable finite flat extension over~\(A\)
  that kills the extension in question.
\item Take the example in which \(A\) is the local ring of an ordinary
  double point of a curve and \(B\) is the local ring at a point of the
  normalization.  Then \(A\subset B\), and \(B\) is unramified
  over\marginpar{\raggedright\footnotesize French source p.~4} \(A\) with
  trivial residue-field extension; the map
  \(\widehat A\to\widehat B\) is surjective but \emph{not injective}.
  We shall therefore strengthen the notion of unramifiedness.
\end{itemize}

\section{Étale morphisms. Étale coverings}
\label{I.4}

We shall take for granted whatever is needed concerning flat morphisms; these
facts will be proved later, if necessary\footnote{Cf.\ Expos\'e~IV.}.

\begin{definition}
\label{I.4.1}
\begin{enumerate}
\item[(a)] Let \(f\colon X\to Y\) be a morphism of finite type.  One says
  that \(f\) is \emph{étale} at~\(x\) if \(f\) is flat at~\(x\) and net
  (equivalently, unramified) at~\(x\).  One says that \(f\) is étale if it
  is so at every point.  One says that \(X\) is étale at~\(x\) over~\(Y\),
  or that it is a \(Y\)-prescheme étale at~\(x\), and so on.
\item[(b)] Let \(f\colon A\to B\) be a local homomorphism.  One says that
  \(f\) is étale, or that \(B\) is étale over~\(A\), if \(B\) is flat and
  unramified over~\(A\).\footnote{Cf.\ the regrets in Expos\'e~III,
  no.~1.2.}
\end{enumerate}
\end{definition}

\begin{proposition}
\label{I.4.2}
For \(B/A\) to be étale, it is necessary and sufficient that
\(\widehat B/\widehat A\) be étale.
\end{proposition}

Indeed, this holds separately for ``unramified'' and for ``flat.''

\begin{corollary}
\label{I.4.3}
Let \(f\colon X\to Y\) be of finite type, and let \(x\in X\).  Whether
\(f\) is étale at~\(x\) depends only on the local homomorphism
\(\mathcal O_{f(x)}\to\mathcal O_x\), and even only on the corresponding
homomorphism between the completions.
\end{corollary}

\begin{corollary}
\label{I.4.4}
Suppose that the residue-field extension
\(\kappa(A)\to\kappa(B)\) is trivial, or that \(\kappa(A)\) is
algebraically closed.  Then \(B/A\) is étale if and only if
\(\widehat A\to\widehat B\) is an isomorphism.
\end{corollary}

This follows by combining flatness with Proposition~\ref{I.3.7}.

\begin{proposition}
\label{I.4.5}
Let \(f\colon X\to Y\) be a morphism of finite type.  Then the set of
points at which it is étale is open.
\end{proposition}

Indeed, this holds separately for ``unramified'' and for ``flat.''

This proposition shows that one may dispense with ``pointwise'' statements
when studying morphisms of finite type that are étale at some point.

\begin{proposition}
\label{I.4.6}
\begin{enumerate}
\item[(i)] An open immersion is étale.
\item[(ii)] The composite of two étale morphisms is étale.
\item[(iii)] A base change of an étale morphism is étale.
\end{enumerate}
\end{proposition}

\marginpar{\raggedright\footnotesize French source p.~5}%
Indeed, \textup{(i)} is trivial, and for \textup{(ii)} and \textup{(iii)}
it is enough to note that the assertions hold for ``unramified'' and for
``flat.''  To tell the truth, there are also corresponding statements for
local homomorphisms (without a finiteness condition); in any event, they will
have to appear in the \emph{multiplodoque}, beginning with the unramified
case.

\begin{corollary}
\label{I.4.7}
The Cartesian product of two étale morphisms is likewise étale.
\end{corollary}

\begin{corollary}
\label{I.4.8}
Let \(X\) and \(X'\) be of finite type over~\(Y\), and let
\(g\colon X\to X'\) be a \(Y\)-morphism.  If \(X'\) is unramified over
\(Y\) and \(X\) is étale over~\(Y\), then \(g\) is étale.
\end{corollary}

Indeed, \(g\) is the composite of the graph morphism
\[
  \Gamma_g\colon X\longrightarrow X\times_Y X',
\]
which is an open immersion by Corollary~\ref{I.3.4}, and the projection
\(X\times_Y X'\to X'\), which is étale because it is obtained from the
étale morphism \(X\to Y\) by the base change \(X'\to Y\).

\begin{definition}
\label{I.4.9}
An \emph{étale covering} (respectively, a \emph{net covering}, also called
an unramified covering) of~\(Y\) is a \(Y\)-scheme \(X\) that is finite
over~\(Y\) and étale (respectively, unramified) over~\(Y\).
\end{definition}

The first condition means that \(X\) is defined by a coherent sheaf of
algebras \(\mathcal B\) on~\(Y\).  The second then means that \(\mathcal B\)
is locally free over~\(Y\) (with no corresponding local-freeness requirement
in the net case), and, moreover, that for every \(y\in Y\) the fiber
\[
  \mathcal B(y)=\mathcal B_y\otimes_{\mathcal O_y}\kappa(y)
\]
is a separable algebra over~\(\kappa(y)\), that is, a finite product of
finite separable field extensions of~\(\kappa(y)\).

\begin{proposition}
\label{I.4.10}
Let \(X\) be a flat covering of~\(Y\) of degree~\(n\) (the definition of
this term deserved to appear in Definition~\ref{I.4.9}), defined by a
coherent locally free sheaf of algebras \(\mathcal B\).  In the usual way one
defines the trace homomorphism \(\mathcal B\to\mathcal A\), which is a
homomorphism of \(\mathcal A\)-modules, where
\(\mathcal A=\mathcal O_Y\).  For \(X\) to be étale, it is necessary and
sufficient that the corresponding trace pairing
\((x,y)\longmapsto\operatorname{Tr}_{\mathcal B/\mathcal A}(xy)\) induce an
isomorphism \(\mathcal B\xrightarrow{\sim}\mathcal B^\vee\),\footnote{The French TeX and the
original printing instead give \(\mathcal B\) as the target.  An
\(\mathcal A\)-bilinear form induces a map
\(\mathcal B\to\mathcal B^\vee\); the dual is therefore the type-correct
target, used here with the source-original defect disclosed.} or,
equivalently,
that the \emph{discriminant section}
\[
  d_{X/Y}=d_{\mathcal B/\mathcal A}
  \in
  \Gamma\!\left(
    Y,
    \bigwedge\nolimits^n\mathcal B^\vee
    \otimes_{\mathcal A}
    \bigwedge\nolimits^n\mathcal B^\vee
  \right)
\]
be invertible, or, finally, that the discriminant ideal defined by this
section be the unit ideal.
\end{proposition}

Indeed, one reduces to the case \(Y=\operatorname{Spec}(k)\); the assertion
is then the well-known separability criterion, and becomes immediate after
passing to an algebraic closure of~\(k\).

\marginpar{\raggedright\footnotesize French source p.~6}%
\paragraph{Remark.}
There will be a less trivial statement below, when one does not assume a
priori that \(X\) is flat over~\(Y\), but instead imposes a normality
hypothesis.

\input{drafts/SGA1_I_5_English_source_draft.texfrag}

\input{drafts/SGA1_I_6_English_source_draft.texfrag}

\end{document}
